\question{25}{Considere un número $x$, inicialmente igual a $1$, sobre el que se puede aplicar repetidamente una de tres operaciones:

\begin{itemize}
  \item Multiplicar $x$ por $2$.
  \item Multiplicar $x$ por $3$.
  \item Sumarle $1$ a $x$.
\end{itemize}

Dado un número natural $n$ ($n \ge 1$), encuentre el mínimo número de operaciones necesarias para obtener el número $n$ empezando desde $1$.

E.g.:

\begin{itemize}
  \item $n=1$: $0$ operaciones ($1$).
  \item $n=5$: $3$ operaciones ($1, 2, 4, 5$).
  \item $n=10$: $3$ operaciones ($1, 3, 9, 10$).
  \item $n=111$: $6$ operaciones ($1, 2, 4, 12, 36, 37, 111$).
\end{itemize}
}

\begin{enumerate}
  \item (10 décimas) Defina una relación de recurrencia que resuelva el problema. Explique en palabras qué significa la función y qué significa cada variable. Sea sucinto.
  \item (5 décimas) Dibuje o describa el grafo de dependencias de la relación de recurrencia. ¿Cuánto espacio requiere como mínimo una solución de programación dinámica para este problema?
  \item (10 décimas) Implemente una solución de programación dinámica, ya sea bottom-up (tabulación) o top-down (memoización), en Python, Java, C, C++, JavaScript o TypeScript.
\end{enumerate}

\bigskip
